%==============================================================================
\section{Scatter Matrix Plot --- Ma trận scatter}
\label{sec:scatter-matrix}
%==============================================================================

\begin{dinhnghia}[title={Scatter Matrix Plot}]
Scatter Matrix Plot (còn gọi \textbf{Pair Plot}) là biểu diễn quan hệ giữa \textbf{nhiều biến} cùng lúc.
Mỗi ô là scatter plot của một \textbf{cặp feature}; đường chéo có thể hiển thị phân phối từng biến.
\end{dinhnghia}

\begin{ynghia}[title={Vai trò trong ML}]
Công cụ hữu ích để trực quan hoá tương quan giữa các feature, hỗ trợ feature engineering và hiểu cấu trúc dữ liệu trước khi huấn luyện.
\end{ynghia}

\subsection{Scatter matrix trong Python}

\begin{phuongphap}[title={Dataset và thư viện}]
Ví dụ dùng \textbf{Iris} qua Seaborn: 4 feature (sepal/petal length/width), 150 mẫu, 3 loài Setosa, Versicolor, Virginica.
Seaborn xây trên Matplotlib, API gọn cho pair plot.
\end{phuongphap}

\subsection{Ví dụ --- pairplot với màu theo loài}

\begin{vidu}[title={Code}]
\begin{lstlisting}
import seaborn as sns
import pandas as pd
import matplotlib.pyplot as plt

# load iris dataset
iris = sns.load_dataset('iris')

# create scatter matrix plot
sns.pairplot(iris, hue='species')

# show plot
plt.show()
\end{lstlisting}

\texttt{load\_dataset('iris')} nạp bảng Iris có cột \texttt{species}.
\texttt{pairplot(..., hue='species')} tô màu điểm theo loài.
\end{vidu}

\begin{ynghia}[title={Đọc biểu đồ}]
Ma trận gồm scatter từng cặp feature; màu theo loài giúp thấy cụm tách biệt (ví dụ Setosa tách rõ ở nhiều cặp biến).
\end{ynghia}

\begin{vidu}[title={Output}]
\tsimg{06_scatter_matrix_plot/img01.jpg}
\end{vidu}

\begin{luuy}[title={Pandas scatter\_matrix}]
Ngoài \texttt{sns.pairplot}, có thể dùng \texttt{pd.plotting.scatter\_matrix} với tham số \texttt{diagonal='hist'}.
\end{luuy}
